% ============================================================
% MÓDULO: GESTIÓN DE LA CALIDAD Y AMBIENTAL. ANÁLISIS FODA
% ============================================================
\clearpage
\section*{MÓDULO: GESTIÓN DE LA CALIDAD Y AMBIENTAL}
\subsection*{Normas ISO 9001, ISO 14001 y análisis FODA}

\begin{recuadro}
\textbf{Objetivo del módulo:} comprender qué es un sistema de gestión de la calidad y un
sistema de gestión ambiental, cómo se relacionan con la actividad aeronáutica y cómo usar el
análisis FODA para diagnosticar una organización o un proyecto.
\end{recuadro}

% ------------------------------------------------------------
\section*{1. CALIDAD Y NORMAS}

\textbf{Calidad} es el grado en que un producto o servicio cumple con los requisitos del
cliente y con las normas aplicables. En aviación, además, la calidad está directamente
ligada a la \textbf{seguridad}.

Una \textbf{norma} es un documento, aprobado por consenso, que establece requisitos para
hacer algo de manera uniforme. La \textbf{ISO} (Organización Internacional de Normalización)
publica normas internacionales; en Argentina el organismo de normalización es el
\textbf{IRAM}, que publica las versiones IRAM-ISO.

% ------------------------------------------------------------
\section*{2. ISO 9001: SISTEMA DE GESTIÓN DE LA CALIDAD}

La \textbf{ISO 9001} establece los requisitos de un \textbf{Sistema de Gestión de la Calidad
(SGC)}. No dice cómo fabricar un producto: dice cómo organizar la empresa para hacer las
cosas bien siempre, satisfacer al cliente y mejorar continuamente. Es certificable.

\subsubsection*{2.1. Principios de la gestión de la calidad}
\begin{enumerate}
  \item Enfoque al cliente.
  \item Liderazgo.
  \item Compromiso de las personas.
  \item Enfoque a procesos.
  \item Mejora continua.
  \item Toma de decisiones basada en evidencia (datos).
  \item Gestión de las relaciones (con proveedores y otras partes interesadas).
\end{enumerate}

\subsubsection*{2.2. Elementos de un SGC}
\begin{itemize}
  \item \textbf{Política y objetivos de calidad} medibles.
  \item \textbf{Mapa de procesos:} procesos estratégicos, operativos (los que generan el
        producto o servicio) y de apoyo (compras, pañol, mantenimiento, RR.\,HH.).
  \item \textbf{Información documentada:} procedimientos, instrucciones de trabajo,
        registros (lo que no está registrado, no se hizo).
  \item \textbf{Gestión de riesgos y oportunidades.}
  \item \textbf{Control de proveedores} y del \textbf{equipamiento de medición} (calibración).
  \item \textbf{No conformidades y acciones correctivas:} cuando algo sale mal, se registra,
        se busca la causa raíz y se corrige para que no se repita.
  \item \textbf{Auditorías internas} y \textbf{revisión por la dirección}.
\end{itemize}

\subsubsection*{2.3. El ciclo PDCA en ISO 9001}
La norma está construida sobre el ciclo \textbf{Planificar -- Hacer -- Verificar -- Actuar},
el mismo de Kaizen: planificar el sistema, operarlo, medir y auditar, y mejorar.

\subsubsection*{2.4. Certificación}
Un \textbf{organismo certificador} independiente audita la empresa; si cumple, emite un
certificado válido por tres años, con auditorías de seguimiento anuales.

\paragraph*{Calidad en la aviación.}
En la industria aeronáutica, sobre la base de ISO 9001 se usan requisitos más exigentes:
\begin{itemize}
  \item \textbf{AS9100 / EN 9100:} normas de calidad para fabricantes del sector aeroespacial
        y de defensa (agregan trazabilidad, gestión de la configuración, prevención de piezas
        falsificadas y de FOD, factores humanos).
  \item \textbf{RAAC 145:} los talleres aeronáuticos de reparación deben tener un
        \textbf{Manual de Procedimientos de Inspección} y un sistema de calidad aprobado por la
        ANAC. La certificación de la autoridad es obligatoria; la ISO es voluntaria.
\end{itemize}

% ------------------------------------------------------------
\section*{3. ISO 14001: SISTEMA DE GESTIÓN AMBIENTAL}

La \textbf{ISO 14001} establece los requisitos de un \textbf{Sistema de Gestión Ambiental
(SGA)}: identificar cómo la actividad de la empresa afecta al ambiente, cumplir la legislación
ambiental y reducir los impactos de forma continua.

\begin{itemize}
  \item \textbf{Aspecto ambiental:} elemento de la actividad que interactúa con el ambiente
        (consumo de energía, generación de residuos, uso de químicos).
  \item \textbf{Impacto ambiental:} el cambio que produce en el ambiente (contaminación del
        suelo o del agua, emisiones).
  \item \textbf{Requisitos legales:} leyes nacionales, provinciales y municipales.
  \item \textbf{Objetivos ambientales} medibles, controles operativos y preparación ante
        emergencias (derrames, incendios).
\end{itemize}

\subsubsection*{3.1. Aspectos ambientales de un taller aeronáutico y electrónico}
\begin{center}
\small
\begin{tabular}{|p{0.30\linewidth}|p{0.60\linewidth}|}
\hline
\textbf{Aspecto} & \textbf{Buena práctica} \\ \hline
Residuos de aparatos eléctricos y electrónicos (RAEE): placas, cables, equipos & Separarlos y
entregarlos a operadores habilitados; no tirarlos con la basura común. \\ \hline
Baterías (plomo-ácido, níquel-cadmio, litio) & Almacenamiento seguro y disposición como
residuo peligroso. \\ \hline
Aceites, combustibles, solventes, pinturas & Bateas antiderrame, kits de derrame,
recipientes rotulados, retiro por transportista habilitado. \\ \hline
Soldadura (humos, estaño con plomo) & Extracción de humos, preferir aleaciones sin plomo cuando
la norma del equipo lo permita. \\ \hline
Energía & Apagar equipos, iluminación LED, mantenimiento de compresores. \\ \hline
\end{tabular}
\end{center}

En Argentina, los \textbf{residuos peligrosos} están regulados por la \textbf{Ley 24.051} y
normas provinciales: quien los genera debe inscribirse, almacenarlos de forma segura y
entregarlos solo a transportistas y operadores habilitados, conservando los
\textbf{manifiestos} (comprobantes) de cada retiro.

% ------------------------------------------------------------
\section*{4. ANÁLISIS FODA}

El \textbf{FODA} (en inglés, SWOT) es una herramienta para diagnosticar la situación de una
empresa, un proyecto o incluso de una persona, y definir estrategias.

\begin{center}
\small
\begin{tabular}{|p{0.13\linewidth}|p{0.37\linewidth}|p{0.37\linewidth}|}
\hline
 & \textbf{Positivo} & \textbf{Negativo} \\ \hline
\textbf{Interno} (lo controlamos) & \textbf{Fortalezas:} lo que hacemos bien, recursos
propios. & \textbf{Debilidades:} lo que nos falta o hacemos mal. \\ \hline
\textbf{Externo} (no lo controlamos) & \textbf{Oportunidades:} tendencias o situaciones del
entorno que podemos aprovechar. & \textbf{Amenazas:} situaciones del entorno que nos pueden
perjudicar. \\ \hline
\end{tabular}
\end{center}

\subsubsection*{4.1. Ejemplo: taller de aviónica recién formado por egresados}
\begin{center}
\small
\begin{tabular}{|p{0.44\linewidth}|p{0.44\linewidth}|}
\hline
\textbf{Fortalezas} & \textbf{Debilidades} \\
\begin{itemize}[leftmargin=*,nosep]
  \item Formación técnica específica y actualizada.
  \item Equipo joven, con manejo de software y redes.
  \item Costos fijos bajos.
\end{itemize}
&
\begin{itemize}[leftmargin=*,nosep]
  \item Poca experiencia y pocos clientes.
  \item Instrumental limitado.
  \item Falta de habilitación ANAC propia al inicio.
\end{itemize}
\\ \hline
\textbf{Oportunidades} & \textbf{Amenazas} \\
\begin{itemize}[leftmargin=*,nosep]
  \item Flota de aviación general antigua que necesita modernizar su aviónica.
  \item Crecimiento de drones y aviación agrícola.
  \item Pocos talleres especializados en la zona.
\end{itemize}
&
\begin{itemize}[leftmargin=*,nosep]
  \item Inestabilidad económica y del tipo de cambio.
  \item Costo de importación de componentes.
  \item Competencia de talleres establecidos.
\end{itemize}
\\ \hline
\end{tabular}
\end{center}

\subsubsection*{4.2. Del diagnóstico a la estrategia (matriz FODA cruzada)}
\begin{itemize}
  \item \textbf{FO (ofensivas):} usar fortalezas para aprovechar oportunidades. \textit{Ofrecer
        servicios para drones y modernización de aviónica con presencia fuerte en redes.}
  \item \textbf{DO (de reorientación):} superar debilidades aprovechando oportunidades.
        \textit{Asociarse con un taller habilitado mientras se obtiene la certificación propia.}
  \item \textbf{FA (defensivas):} usar fortalezas para reducir amenazas. \textit{Mantener costos
        fijos bajos para resistir los ciclos económicos.}
  \item \textbf{DA (de supervivencia):} reducir debilidades y evitar amenazas. \textit{Comprar
        instrumental usado calibrado y armar un stock mínimo de componentes críticos.}
\end{itemize}

% ------------------------------------------------------------
\section*{5. ACTIVIDADES DEL MÓDULO}

\begin{enumerate}
  \item \textbf{Mapa de procesos del taller de la escuela:} identificar procesos estratégicos,
        operativos y de apoyo, y redactar un procedimiento (por ejemplo, préstamo de
        instrumental del pañol) con su registro.
  \item \textbf{No conformidad:} a partir de un error real o simulado del taller, completar un
        informe de no conformidad con causa raíz (5 porqués) y acción correctiva.
  \item \textbf{Relevamiento ambiental:} listar los aspectos ambientales del taller y proponer
        un plan de gestión de RAEE y baterías.
  \item \textbf{FODA:} realizar el FODA y la matriz cruzada del proyecto final del grupo y un
        FODA personal como futuro técnico.
\end{enumerate}
